\section{Survival at equal quality}\label{sec:survival}

\paragraph{Main estimates.} Table~\ref{tab:survival} reports $\delta$ in equation~\eqref{eq:did}, the extra
domestic--foreign gap in emerging-market offices. Across the ten offices and for all firms and universities together,
without quality controls it is $-10.1$ points (column 1). Comparing patents with the same citations, technology field
and owner group, it becomes $+4.1$ (column 2): the raw shortfall reflects that domestic patents in emerging-market
offices are much less cited than foreign ones. Columns 3--6 estimate it separately for each owner group. Measured against foreign owners in the same office,
emerging-market firms of every size keep their patents longer than advanced-economy firms do: small firms by 2.8 points,
mid-sized firms by 7.9 and large firms by 9.0. Universities are behind, at $-5.7$, though not significantly.

\paragraph{Which offices drive it.} China and the United States are the largest offices on each side, so
the lower two blocks of Table~\ref{tab:survival} separate them. China against the
United States alone gives the same pooled estimate, $+3.9$, with Chinese small firms 5.4 points ahead and
Chinese universities 15.3 points behind. Without China and the United States the pooled estimate is close to
zero, $-1.0$; mid-sized and large firms are still ahead, by 2.4 and 3.5 points, the second not
significantly, but small firms in Korea, Russia, Taiwan
and the Czech Republic are 4.4 points behind small firms in Japan, Germany, the United Kingdom and France. The
benchmark matters here: small Japanese firms' patents are 18 points more likely to be alive than foreign owners' patents
in Japan, by far the largest domestic advantage among the advanced-economy offices (Supplemental Appendix Table~S4), and
without the United States Japan sets the bar.

\begin{table}[t]
\centering
\caption{Extra domestic--foreign survival gap in emerging-market offices}\label{tab:survival}
\small
% Generated by em-patent-value code/38_paper_exhibits.py (main_table2.csv, step 37); do not edit by hand
\begin{tabular*}{\linewidth}{@{\extracolsep{\fill}}lcccccc}
\toprule
Alive 8 years after grant & \multicolumn{2}{c}{All owners} & Small & Mid-sized & Large & \\
 & & & firms & firms & firms & Universities \\
\cmidrule(lr){2-3}
 & (1) & (2) & (3) & (4) & (5) & (6) \\
\midrule
\multicolumn{7}{l}{\emph{All ten offices}} \\
Domestic owner $\times$ EM office & $-12.3^{***}$ & $3.5^{*}$ & $0.5$ & $8.0^{***}$ & $5.4^{**}$ & $-14.2^{***}$ \\
 & (1.7) & (1.8) & (0.7) & (1.3) & (2.5) & (2.9) \\
Domestic owner & $3.8^{***}$ & $5.0^{***}$ & $3.8^{***}$ & $5.4^{***}$ & $5.0^{***}$ & $4.9^{**}$ \\
 & (0.9) & (1.0) & (0.3) & (0.4) & (1.3) & (2.2) \\
Observations & 2{,}504{,}283 & 2{,}504{,}279 & 255{,}234 & 288{,}177 & 1{,}808{,}726 & 152{,}115 \\
\addlinespace
\multicolumn{7}{l}{\emph{Excl. China and US}} \\
Domestic owner $\times$ EM office & $-9.8^{***}$ & $0.9$ & $-4.8^{***}$ & $5.1^{**}$ & $1.7$ & $-14.4^{**}$ \\
 & (2.2) & (2.4) & (1.5) & (2.0) & (3.1) & (6.1) \\
Domestic owner & $4.9^{***}$ & $8.8^{***}$ & $8.0^{***}$ & $8.1^{***}$ & $9.0^{***}$ & $16.8^{***}$ \\
 & (1.2) & (1.2) & (0.8) & (1.0) & (1.6) & (4.7) \\
Observations & 1{,}229{,}431 & 1{,}229{,}427 & 101{,}703 & 126{,}004 & 938{,}771 & 62{,}922 \\
\addlinespace
\multicolumn{7}{l}{\emph{China and US only}} \\
Domestic owner $\times$ EM office & $-15.3^{***}$ & $3.6$ & $2.6^{***}$ & $10.3^{***}$ & $6.8$ & $-19.7^{***}$ \\
 & (2.6) & (2.8) & (0.8) & (1.7) & (4.4) & (2.8) \\
Domestic owner & $3.4^{***}$ & $3.8^{***}$ & $2.3^{***}$ & $3.6^{***}$ & $3.4^{*}$ & $4.7^{**}$ \\
 & (1.0) & (1.4) & (0.4) & (0.5) & (1.9) & (2.1) \\
Observations & 1{,}274{,}852 & 1{,}274{,}852 & 153{,}531 & 162{,}173 & 869{,}955 & 89{,}193 \\
\midrule
\multicolumn{7}{l}{\emph{Controls (by office)}} \\
\quad Filing year & \checkmark & \checkmark & \checkmark & \checkmark & \checkmark & \checkmark \\
\quad Field, citation bin, grant lag & & \checkmark & \checkmark & \checkmark & \checkmark & \checkmark \\
\quad Domestic $\times$ grant lag & & \checkmark & \checkmark & \checkmark & \checkmark & \checkmark \\
\quad Owner group & & \checkmark & -- & -- & -- & -- \\
\bottomrule
\end{tabular*}

\vspace{4pt}
\begin{minipage}{\linewidth}\footnotesize\emph{Notes:} Linear probability models of survival eight years after grant; granted patents of firms and universities filed 2000--2005, grant lag up to ten years; coefficients in percentage points. The coefficient on domestic owner $\times$ EM office is $\delta$ in equation~\eqref{eq:did}: the extra domestic--foreign gap in the emerging-market offices (China, Korea, Russia, Taiwan, Czech Republic) relative to the advanced-economy offices (United States, Japan, Germany, United Kingdom, France). Column 1: office $\times$ filing-year effects only. Columns 2--6: office-specific effects of filing year, technology field (35), citation bin (forward citations to the family: 0, 1, 2, 3--4, 5--9, 10--19, 20 or more) and grant lag in whole years, a common domestic $\times$ grant-lag slope, and in column 2 office $\times$ owner-group effects; columns 3--6 are estimated separately for each owner group. The domestic-owner coefficient is the advanced-economy gap at the mean grant lag. Firm size: inventions filed worldwide in 2000--2005 by the owner. Standard errors clustered by owner in parentheses; inference across offices in \SuppTab{S3}{tab:rich}. $^{*}$~$p<0.1$, $^{**}$~$p<0.05$, $^{***}$~$p<0.01$.
\end{minipage}
\end{table}


\paragraph{Richer quality controls.} Citations are a noisy measure of quality. Adding citations that block later
claims, citations received in the four years after filing and office-specific slopes in the grant lag lowers the estimates for every
firm size by 1--3 points (Supplemental Appendix Table~S3). Small emerging-market firms then come out level ($-0.4$), or 1.7 points behind when family size is also held
fixed, although family size itself partly reflects an owner's capacity to file abroad;
mid-sized and large firms remain 4--8 points ahead.

\clearpage
\newpage 
\paragraph{By country.} Supplemental Appendix Table~S2 compares each emerging-market office with the five advanced-economy
offices, with the specifications of Table~\ref{tab:survival}. Chinese and Korean firms of every size are ahead: small firms by 4.0 and 5.1
points, larger firms by 7--11 points. Small Czech firms are ahead by a similar amount, 6.3 points, though not
significantly. Russia is the exception: small Russian firms are 4.3 points behind, and Russian universities 11.9. Small
Taiwanese firms are 2.4 points behind, while larger Taiwanese firms and universities are ahead. For all owners
together, Korea and Taiwan are ahead, China and the Czech Republic level and Russia behind; in China, domestic firms'
advantage is offset by universities, 16.9 points behind.
Within each office, the domestic--foreign gap for small firms is close to zero in Russia and Taiwan, against about 2--18
points in the advanced-economy offices (Supplemental Appendix Table~S4).
